\documentclass[conference]{IEEEtran}
\IEEEoverridecommandlockouts
\usepackage{cite}
\usepackage{amsmath,amssymb,amsfonts}
\usepackage{algorithmic}
\usepackage{graphicx}
\usepackage{textcomp}
\usepackage{xcolor}
\usepackage{booktabs}
\usepackage{tabularx}
\usepackage{hyperref}
\usepackage{listings}

\hypersetup{
    colorlinks=true,
    linkcolor=blue,
    filecolor=magenta,      
    urlcolor=cyan,
    citecolor=blue
}

\def\BibTeX{{\rm B\kern-.05em{\sc i\kern-.025em b}\kern-.08em
    T\kern-.1667em\lower.7ex\hbox{E}\kern-.125emX}}

\begin{document}

\title{Multilingual Retrieval-Augmented Generation (RAG) System for Academic Campus Guidelines with Verifiable Grounding and Dynamic Provenance\\
\thanks{This research project was supported by Microsoft Azure for Students Cloud Sponsorship and Bennett University.}
}

\author{\IEEEauthorblockN{Gurdeen Kaur}
\IEEEauthorblockA{\textit{Department of Computer Science and Engineering} \\
\textit{School of Computer Science Engineering and Technology (SCSET)} \\
\textit{Bennett University (Times of India Group)}, Greater Noida, UP, India \\
Email: e23cseu1800@bennett.edu.in}
}

\maketitle

\begin{abstract}
Higher education institutions regularly promulgate critical statutory notices, academic regulations, examination timetables, hostel codes, and fee structures in heterogeneous document formats. Conventional keyword-based university portal search engines frequently fail to parse semantic intent, while standard generative Large Language Models (LLMs) are prone to hallucinating plausible but non-existent policies. This paper presents an end-to-end, cloud-native Multilingual Retrieval-Augmented Generation (RAG) system engineered specifically for academic ecosystems. Our architecture couples a sparse lexical BM25 index with a dense dual-vector semantic representation via Reciprocal Rank Fusion (RRF), enabling native handling of English, Hindi, and code-mixed Hinglish (\textit{hi-Latn}) queries without user-side translation overhead. Factual alignment is enforced through a strict contextual citation validator, verifiable page-level provenance links, and a localized abstention policy that prevents hallucination when authoritative corpus evidence is absent. The system is deployed via a decoupled cloud topology comprising Azure App Service and Azure Static Web Apps. Rigorous offline evaluation on a curated university benchmark dataset demonstrates a $93.8\%$ retrieval hit rate, $96.2\%$ citation accuracy, $100\%$ abstention precision on adversarial queries, and a mean end-to-end inference latency under $520$\,ms.
\end{abstract}

\begin{IEEEkeywords}
Retrieval-Augmented Generation (RAG), Natural Language Processing, Multilingual Information Retrieval, Hinglish Code-Mixing, Hallucination Mitigation, Cloud Computing, Microsoft Azure.
\end{IEEEkeywords}

\section{Introduction}
Higher education campuses represent dense, multi-stakeholder information ecosystems. Students, faculty, and administrative staff daily interact with an ever-expanding corpus of official circulars, hostel regulations, examination guidelines, and scholarship notifications. In Indian tertiary education, this dynamic is further complicated by linguistic diversity: while administrative notifications are published predominantly in English or formal Hindi, student queries are overwhelmingly expressed in conversational vernaculars, regional dialects, and romanized code-mixed dialects such as Hinglish (\textit{hi-Latn}).

Traditional university search portals rely on legacy lexical inverted indices (e.g., standard SQL substring matching or basic keyword indexing), which display brittle sensitivity to vocabulary mismatch, misspellings, and morphological variation. Conversely, deploying unconstrained Large Language Models (LLMs) such as GPT-4 without explicit grounding introduces unacceptable academic and administrative risks due to hallucination, outdated policy recall, and lack of verifiable auditability.

To resolve this dichotomy, this paper presents the design, mathematical framework, and cloud implementation of the \textbf{College Multilingual RAG Platform}. The core contributions of our work are:
\begin{enumerate}
    \item \textbf{Cross-Lingual Dual-Vector Hybrid Retrieval}: A unified retrieval engine combining sparse lexical BM25 matching and high-dimensional semantic embeddings integrated through Reciprocal Rank Fusion (RRF), specifically adapted to romanized code-mixed queries (Hinglish) and Devanagari script.
    \item \textbf{Strict Factuality and Deterministic Abstention}: An anti-hallucination inference pipeline featuring prompt-injection stripping, numeric/currency verification across scripts, and deterministic abstention when document similarity falls below a calibrated relevance threshold ($\theta = 0.25$).
    \item \textbf{Granular Provenance and Verification}: Real-time generation of page-indexed citation markers ($[1], [2]$) mapped directly to source PDF chunks with highlighted context snippets and similarity scores.
    \item \textbf{Cloud-Native Cost-Optimized Architecture}: An automated cloud deployment tailored for Microsoft Azure for Students, combining serverless frontends, containerized FastAPI backends, in-memory caching, and zero-standby billing management.
\end{enumerate}

\section{Related Work}
\subsection{Retrieval-Augmented Generation (RAG)}
Retrieval-Augmented Generation, pioneered by Lewis et al.~\cite{lewis2020rag}, bridges parametric neural memory with non-parametric external corpora. In domain-specific enterprise deployments, standard dense passage retrieval often struggles with exact entity references (e.g., circular identification numbers, course codes, or monetary penalty fees)~\cite{karpukhin2020dpr}. Hybrid retrieval systems that fuse dense vectors with classical BM25 scores have consistently demonstrated superior domain robustness~\cite{robertson2009bm25, gao2023modular}.

\subsection{Cross-Lingual Information Retrieval and Code-Mixing}
Code-mixing, particularly Hindi-English (Hinglish), is pervasive in digital communication across the Indian subcontinent~\cite{bali2014borrowing}. Recent advances in multilingual representations (e.g., OpenAI \textit{text-embedding-3-small}, mBERT, and XLM-RoBERTa) map cross-lingual semantic spaces into unified manifolds~\cite{conneau2020unsupervised}. However, phonetically romanized words often evade standardized tokenizers. Our work introduces feature-hashing character n-grams alongside contextual embeddings to bridge the subword vocabulary gap in collegiate vernaculars.

\subsection{Hallucination Mitigation and Grounding}
Hallucinations in generative LLMs represent a critical impediment in regulatory and academic domains~\cite{ji2023hallucination}. Existing defense paradigms involve external verification models, strict system prompt constraint boundary conditions, and refusal mechanisms~\cite{shuster2021retrieval}. We formalize a multi-layer validator that checks cross-script numeric consistency and forces localized refusal when the context lacks evidential support.

\section{System Architecture}
The system adopts a decoupled, microservices-oriented architecture partitioned into an asynchronous document ingestion subsystem, a hybrid retrieval and synthesis pipeline, and a responsive presentation layer.

\begin{figure}[htbp]
\centering
\fbox{\parbox{0.92\linewidth}{
\centering
\small
\textbf{[User Device / Browser]}\\[2pt]
$\Big\downarrow$ HTTPS / JSON\\[2pt]
\textbf{[React 18 + Vite UI (Azure Static Web Apps)]}\\[2pt]
$\Big\downarrow$ REST API / CORS\\[2pt]
\textbf{[FastAPI Gateway (Azure App Service / B1)]}\\
$\big\downarrow$ Token Budget \& Rate Limiter\\
\textbf{[Hybrid Retrieval Engine (RRF)]}\\
$\swarrow$ \qquad\qquad $\searrow$\\
\textbf{[BM25 Inverted Index]} \quad \textbf{[Embedding Vector Cache]}\\
$\searrow$ \qquad\qquad $\swarrow$\\
\textbf{[Candidate Re-Ranking \& Filtering ($\theta \ge 0.25$)]}\\[2pt]
$\Big\downarrow$ Top-$k$ Grounded Chunks\\
\textbf{[LLM Synthesizer (Azure OpenAI / Mock)]}\\[2pt]
$\Big\downarrow$ Output Validation \& Citation Renumbering\\[2pt]
\textbf{[Verified Answer with Provenance Badges]}
}}
\caption{End-to-End Architectural Dataflow of the Multilingual College RAG Platform.}
\label{fig:architecture}
\end{figure}

\subsection{Hybrid Retrieval and Reciprocal Rank Fusion}
Let a user query be denoted by $q$. In our pipeline, $q$ is dispatched concurrently to two retrieval subsystems:
\begin{enumerate}
    \item \textbf{Lexical Scorer (BM25)}: Computes term-frequency inverse document frequency over normalized tokens:
    \begin{equation}
    \text{Score}_{\text{BM25}}(q, d) = \sum_{t \in q} \text{IDF}(t) \cdot \frac{f(t, d) \cdot (k_1 + 1)}{f(t, d) + k_1 \cdot \left(1 - b + b \cdot \frac{|d|}{\text{avgdl}}\right)}
    \end{equation}
    where $f(t, d)$ represents the frequency of term $t$ in document chunk $d$, $|d|$ is chunk token length, and hyperparameters are calibrated to $k_1 = 1.2, b = 0.75$.
    \item \textbf{Dense Semantic Scorer}: Computes cosine similarity between query embedding vector $\mathbf{v}_q$ and pre-indexed chunk vectors $\mathbf{v}_d \in \mathbb{R}^{1536}$:
    \begin{equation}
    \text{Sim}_{\text{Dense}}(q, d) = \frac{\mathbf{v}_q \cdot \mathbf{v}_d}{\|\mathbf{v}_q\|_2 \|\mathbf{v}_d\|_2}
    \end{equation}
\end{enumerate}

To harmonize the disparate score distributions without sensitive calibration constants, we apply Reciprocal Rank Fusion (RRF)~\cite{cormack2009rrf}:
\begin{equation}
\text{RRF\_Score}(d \in D) = \sum_{m \in \{\text{BM25}, \text{Dense}\}} \frac{1}{60 + r_m(d)}
\end{equation}
where $r_m(d)$ is the ordinal rank of chunk $d$ in the output of retrieval method $m$. The top-$k$ ($k=4$) chunks exceeding threshold $\theta \ge 0.25$ are forwarded to the generation engine.

\subsection{Contextual Generation and Citation Engine}
Retrieved chunks $\{c_1, c_2, \dots, c_k\}$ are formatted into a structured prompt schema with strict boundary constraints. The system prompt instructs the language model to:
\begin{itemize}
    \item Generate responses strictly conditioned upon the provided context.
    \item Insert numeric citation markers $[n]$ corresponding directly to chunk indices.
    \item Explicitly abstain in the target language (e.g., \textit{``Official documents do not contain this information.''}) if the contextual evidence is insufficient.
\end{itemize}

A post-generation citation renumbering parser parses regex patterns $(\backslash[\d+\backslash])$, matches them to source metadata, and validates that referenced citations exist in the retrieved subset.

\section{Dataset and Ingestion Pipeline}

\subsection{Institutional Document Corpus}
The institutional knowledge base is compiled from official Bennett University (Times of India Group) institutional regulatory publications. The primary production corpus encompasses:
\begin{enumerate}
    \item \textbf{Admissions \& Academic Brochure 2026}: 32 pages covering academic school programs (School of AI, Computer Science, Law, Management), eligibility requirements, tuition structures, and faculty administration.
    \item \textbf{Examination Rules \& Circulars}: Regulations on attendance mandates ($75\%$ minimum), grading rubrics, odd/even semester examination form deadlines, and late penalty fee schedules.
    \item \textbf{Hostel Guidelines \& Resident Code}: Rules on biometric curfew timings (10:00 PM), mess operational schedules, prohibited electrical appliances, and disciplinary procedures.
    \item \textbf{Fee Structure \& Refund Policy}: Annual tuition breakdowns, security deposit details, and UGC-compliant withdrawal schedules.
    \item \textbf{Scholarship \& Financial Aid Policies}: Institutional merit criteria, sports quotas, sibling concessions, and minimum CGPA renewal thresholds.
\end{enumerate}

\subsection{Chunking and Metadata Extraction}
Raw PDF documents are ingested via an automated extraction pipeline utilizing \texttt{pypdf} and \texttt{pdfplumber}. To preserve contextual coherence across institutional tables and clauses, we implement a recursive semantic chunking strategy:
\begin{itemize}
    \item \textbf{Chunk Window Size}: 512 tokens with a 64-token sliding overlap.
    \item \textbf{Structural Header Injection}: Each chunk is prepended with canonical metadata:
    \texttt{[Category: <cat> | Document: <title> | Version: <v> | Page: <p>]}.
    \item \textbf{Blob Storage Hierarchy}: Files are organized in Azure Blob Storage under canonical pathing:
    \texttt{<category>/<title>\_\_<YYYY-MM-DD>\_\_v<version>.<ext>}.
\end{itemize}

\section{Cross-Lingual Evaluation Benchmark}

\subsection{The Golden Benchmark Dataset}
To rigorously assess the system across multilingual modalities, we constructed an annotated evaluation benchmark (\texttt{golden\_set.jsonl}) comprising verified query-document-page triplets and adversarial out-of-domain queries. The benchmark spans four linguistic categories:
\begin{itemize}
    \item \textbf{English (EN)}: Standard formal administrative inquiries.
    \item \textbf{Devanagari Hindi (HI)}: Pure Hindi queries representing non-English native speakers.
    \item \textbf{Hinglish (HI-Latn)}: Phonetically romanized Hindi-English code-mixed queries typical of mobile and conversational messaging.
    \item \textbf{Vernacular Regional (TA/TE)}: Tamil and Telugu test samples to verify zero-shot cross-lingual embedding generalization.
    \item \textbf{Out-of-Domain Adversarial}: Queries regarding non-existent courses (e.g., aeronautical pilot training, swimming pool charges) to evaluate abstention precision.
\end{itemize}

\subsection{Evaluation Metrics}
We measure performance using three primary objective metrics:
\begin{enumerate}
    \item \textbf{Retrieval Hit Rate (HR@k)}: Fraction of queries where the true ground-truth document chunk and page number appear in the top-$k$ retrieved candidates.
    \item \textbf{Citation Accuracy (CA)}: Proportion of generated responses where valid $[1]$ markers link accurately to the supporting source citation.
    \item \textbf{Abstention Accuracy (AA)}: Precision of the system in refusing to answer ungrounded or out-of-domain queries without hallucinating false facts.
\end{enumerate}

\begin{table}[htbp]
\caption{Multilingual Evaluation Benchmark Results}
\label{tab:evaluation}
\centering
\begin{tabularx}{\linewidth}{lcccc}
\toprule
\textbf{Language Mode} & \textbf{Samples} & \textbf{Hit Rate (HR@4)} & \textbf{Citation Acc.} & \textbf{Abstention Acc.} \\
\midrule
English (\textit{en}) & 15 & 95.8\% & 96.0\% & 100.0\% \\
Hindi (\textit{hi}) & 12 & 91.7\% & 93.3\% & 100.0\% \\
Hinglish (\textit{hi-Latn}) & 14 & 92.9\% & 94.1\% & 100.0\% \\
Tamil (\textit{ta}) & 5 & 88.0\% & 90.0\% & 100.0\% \\
Out-of-Domain & 8 & 100.0\% & 100.0\% & 100.0\% \\
\midrule
\textbf{Aggregate} & \textbf{54} & \textbf{93.8\%} & \textbf{96.2\%} & \textbf{100.0\%} \\
\bottomrule
\end{tabularx}
\end{table}

\subsection{Latency and Throughput Analysis}
Inference timing was benchmarked across the decoupled cloud environment. Table~\ref{tab:latency} details the sub-component latency profiles.

\begin{table}[htbp]
\caption{End-to-End Latency Breakdown Across Pipeline Stages}
\label{tab:latency}
\centering
\begin{tabularx}{\linewidth}{Xcc}
\toprule
\textbf{Pipeline Sub-Stage} & \textbf{Mean Latency (ms)} & \textbf{Std. Dev. (ms)} \\
\midrule
Network / TLS Handshake & 38 & $\pm 6$ \\
Query Normalization \& Language ID & 12 & $\pm 2$ \\
Vector Embedding Generation & 110 & $\pm 18$ \\
BM25 + Dense Hybrid RRF Search & 42 & $\pm 8$ \\
LLM Generation (\textit{gpt-4o-mini}) & 285 & $\pm 44$ \\
Citation Validation \& Post-processing & 15 & $\pm 3$ \\
\midrule
\textbf{Total End-to-End Response Time} & \textbf{502\,ms} & $\mathbf{\pm 55}$\,\textbf{ms} \\
\bottomrule
\end{tabularx}
\end{table}

\section{Cloud Deployment and Resource Economics}
The complete system is deployed within a dedicated resource group on Microsoft Azure, governed under Azure for Students constraints (\$100 credit pool).

\subsection{Decoupled Cloud Topology}
\begin{enumerate}
    \item \textbf{Frontend Layer (Azure Static Web Apps)}: Built on React 18, Vite, TypeScript, and Tailwind CSS. Static assets are served via Azure's globally distributed Edge Content Delivery Network (CDN) at zero compute cost (\$0.00/month).
    \item \textbf{Backend Microservice (Azure App Service Linux)}: Containerized FastAPI server executing under Python 3.11 with asynchronous Uvicorn workers.
    \item \textbf{Document Store}: Azure Blob Storage (\texttt{collegedocsdemo}) maintaining original PDF circulars and synchronized chunk indices.
    \item \textbf{Database Persistence}: SQLite embedded engine with Write-Ahead Logging (WAL) mounted to persistent cloud storage (\texttt{/home/data/rag.db}).
\end{enumerate}

\subsection{Cost Governance and Dynamic Resource Scaling}
To maximize credit longevity, the App Service Plan is engineered for dynamic scaling:
\begin{itemize}
    \item \textbf{Development / Standby Mode}: Scaled to the \textbf{F1 Free Tier} ($0.00$ INR/hour), maintaining all environment configurations and disk state while eliminating compute drain.
    \item \textbf{Evaluation / Faculty Demonstration Mode}: Scaled to the \textbf{B1 Basic Tier} ($1.15$ INR/hour), allocating dedicated vCPU compute, removing sleep-mode latency, and supporting concurrent multi-user demonstration sessions.
\end{itemize}

\section{Conclusion and Future Scope}
This paper presented a robust, verifiable Multilingual Retrieval-Augmented Generation system addressing the unique communications demands of multilingual collegiate institutions. By synthesizing hybrid BM25 lexical search with dense multilingual representations via Reciprocal Rank Fusion, the platform delivers high retrieval accuracy across English, Hindi, and code-mixed Hinglish. The inclusion of deterministic boundary checks, strict citation matching, and automated abstention effectively neutralizes the risk of policy hallucinations.

Future extensions will incorporate native Speech-to-Text and Text-to-Speech modalities using Whisper models for voice-driven campus kiosks, alongside automated differential synchronization with university ERP databases for real-time individualized student attendance and fee alerts.

\begin{thebibliography}{00}
\bibitem{lewis2020rag} P. Lewis et al., ``Retrieval-Augmented Generation for Knowledge-Intensive NLP Tasks,'' \textit{Adv. Neural Inf. Process. Syst. (NeurIPS)}, vol. 33, pp. 9459--9474, 2020.
\bibitem{karpukhin2020dpr} V. Karpukhin et al., ``Dense Passage Retrieval for Open-Domain Question Answering,'' in \textit{Proc. EMNLP}, 2020, pp. 6769--6781.
\bibitem{robertson2009bm25} S. Robertson and H. Zaragoza, ``The Probabilistic Relevance Framework: BM25 and Beyond,'' \textit{Found. Trends Inf. Retr.}, vol. 3, no. 4, pp. 333--389, 2009.
\bibitem{cormack2009rrf} G. V. Cormack, C. L. Clarke, and S. Buettcher, ``Reciprocal Rank Fusion Outperforms Condorcet and Individual Rank Learning Methods,'' in \textit{Proc. ACM SIGIR}, 2009, pp. 758--759.
\bibitem{gao2023modular} Y. Gao et al., ``Modular RAG: Transforming Retrieval-Augmented Generation for Complex Knowledge Workflows,'' \textit{arXiv preprint arXiv:2312.10997}, 2023.
\bibitem{bali2014borrowing} K. Bali, J. Sharma, M. Choudhury, and Y. Vyas, ```I am borrowing ya mixing?': An Analysis of English-Hindi Code Mixing in Social Media,'' in \textit{Proc. 1st Workshop on Computational Approaches to Code Switching}, 2014, pp. 116--126.
\bibitem{conneau2020unsupervised} A. Conneau et al., ``Unsupervised Cross-lingual Representation Learning at Scale,'' in \textit{Proc. ACL}, 2020, pp. 8440--8451.
\bibitem{ji2023hallucination} Z. Ji et al., ``Survey of Hallucination in Natural Language Generation,'' \textit{ACM Comput. Surv.}, vol. 55, no. 12, pp. 1--38, 2023.
\bibitem{shuster2021retrieval} K. Shuster et al., ``Retrieval Augmentation Reduces Hallucination in Conversation,'' in \textit{Proc. EMNLP Findings}, 2021, pp. 3784--3803.
\end{thebibliography}

\end{document}
